% Options for packages loaded elsewhere
\PassOptionsToPackage{unicode}{hyperref}
\PassOptionsToPackage{hyphens}{url}
\PassOptionsToPackage{dvipsnames,svgnames,x11names}{xcolor}
\documentclass[
  10pt,
]{article}
\usepackage{xcolor}
\usepackage[margin=0.75in]{geometry}
\usepackage{amsmath,amssymb}
\setcounter{secnumdepth}{-\maxdimen} % remove section numbering
\usepackage{iftex}
\ifPDFTeX
  \usepackage[T1]{fontenc}
  \usepackage[utf8]{inputenc}
  \usepackage{textcomp} % provide euro and other symbols
\else % if luatex or xetex
  \usepackage{unicode-math} % this also loads fontspec
  \defaultfontfeatures{Scale=MatchLowercase}
  \defaultfontfeatures[\rmfamily]{Ligatures=TeX,Scale=1}
\fi
\usepackage{lmodern}
\ifPDFTeX\else
  % xetex/luatex font selection
\fi
% Use upquote if available, for straight quotes in verbatim environments
\IfFileExists{upquote.sty}{\usepackage{upquote}}{}
\IfFileExists{microtype.sty}{% use microtype if available
  \usepackage[]{microtype}
  \UseMicrotypeSet[protrusion]{basicmath} % disable protrusion for tt fonts
}{}
\makeatletter
\@ifundefined{KOMAClassName}{% if non-KOMA class
  \IfFileExists{parskip.sty}{%
    \usepackage{parskip}
  }{% else
    \setlength{\parindent}{0pt}
    \setlength{\parskip}{6pt plus 2pt minus 1pt}}
}{% if KOMA class
  \KOMAoptions{parskip=half}}
\makeatother
\usepackage{longtable,booktabs,array}
\usepackage{calc} % for calculating minipage widths
% Correct order of tables after \paragraph or \subparagraph
\usepackage{etoolbox}
\makeatletter
\patchcmd\longtable{\par}{\if@noskipsec\mbox{}\fi\par}{}{}
\makeatother
% Allow footnotes in longtable head/foot
\IfFileExists{footnotehyper.sty}{\usepackage{footnotehyper}}{\usepackage{footnote}}
\makesavenoteenv{longtable}
\setlength{\emergencystretch}{3em} % prevent overfull lines
\providecommand{\tightlist}{%
  \setlength{\itemsep}{0pt}\setlength{\parskip}{0pt}}
\usepackage{enumitem}
\setlist{nosep}
\usepackage{fvextra}
\fvset{fontsize=\footnotesize}
\RecustomVerbatimEnvironment{verbatim}{Verbatim}{fontsize=\footnotesize}
\usepackage{bookmark}
\IfFileExists{xurl.sty}{\usepackage{xurl}}{} % add URL line breaks if available
\urlstyle{same}
\hypersetup{
  pdftitle={Prefill once{,} reuse anywhere: non-prefix KV reuse for agent memory in vLLM},
  pdfauthor={Atharva Mohite, Teppei Kawashima, Prakhar Langer, Vishal Puneyani},
  colorlinks=true,
  linkcolor={Maroon},
  filecolor={Maroon},
  citecolor={Blue},
  urlcolor={Blue},
  pdfcreator={LaTeX via pandoc}}

\title{Prefill once, reuse anywhere: non-prefix KV reuse for agent
memory in vLLM}
\usepackage{etoolbox}
\makeatletter
\providecommand{\subtitle}[1]{% add subtitle to \maketitle
  \apptocmd{\@title}{\par {\large #1 \par}}{}{}
}
\makeatother
\subtitle{CS 6220 Course Project Proposal, Fall 2026, Group 7}
\author{Atharva Mohite, Teppei Kawashima, Prakhar Langer, Vishal
Puneyani}
\date{}

\begin{document}
\maketitle

\section{1. Motivation and objectives}\label{motivation-and-objectives}

Before an LLM produces its first token it runs \emph{prefill}, one pass
over the prompt that computes a key and a value vector for every token
at every layer (the KV cache). Prefill cost grows with prompt length and
dominates time to first token (TTFT) for long prompts.

Memory-augmented agents produce exactly these prompts. For each question
the agent retrieves relevant entries from its long-term memory and
places them ahead of the question. Over a long conversation the same
entries come back many times, in different combinations and orders, and
the serving engine prefills them from scratch each time. vLLM {[}1{]}
reuses KV only through \emph{prefix caching}, which needs two prompts to
start with identical tokens. If one question retrieves memories A, B, C
and the next retrieves C, A, D, the prompts diverge right after the
system prompt, and A and C are prefilled again.

Reusing a chunk's KV at a new position is not exact. A token's KV
depends on its position, which rotary position embeddings (RoPE) {[}2{]}
let us correct exactly by re-rotating the cached keys. It also depends
on the tokens before it, which no rotation fixes: a chunk cached on its
own never attended to the memories now placed in front of it. That
second error is usually small, because tokens attend mostly within their
own chunk and the question is always computed fresh. EPIC {[}3{]},
CacheBlend {[}4{]} and AgentKVShift {[}5{]} close the gap by recomputing
a small subset of tokens, but each assumes the error lives in a
different place, and they have not been compared on one model and one
agent-memory workload. vLLM documents a KV-connector interface, and
LMCache's {[}6{]} development branch describes a CacheBlend operator,
but we found no stable, documented path that runs corrected non-prefix
reuse on the single vLLM release we pin (0.27.1, the version in
LMCache 0.5.4's release image). Two limitations were still open
when we checked: LMCache's blending example is reported out of date for
recent vLLM (issue \#4476), and its newer blend mode has no vLLM-side
client in multi-process mode (issue \#5101).

\textbf{Research question.} \emph{How much prefill can non-prefix KV
reuse remove from a memory-augmented agent, and how do different ways of
correcting reused KV trade answer quality against prefill cost?}

\textbf{Secondary question.} \emph{Served through an unmodified vLLM,
does reuse lower TTFT and raise throughput, and how often must memories
repeat before it beats prefix caching?}

\textbf{Objectives.} We will deliver (1) a Transformers implementation
of chunk-level KV reuse with RoPE re-rotation and four correction
methods behind one shared recompute loop; (2) a controlled comparison of
the four methods on LoCoMo {[}7{]}, measuring F1 against recompute
budget and prefill cost against the number of retrieved memories; and
(3) a vLLM connector plugin that serves reused memory KV to an
unmodified vLLM, with latency and throughput measured under a workload
whose repetition rate we control. Together these tell someone serving
agents or RAG which correction is worth its cost, and how often memories
must repeat before non-prefix reuse beats the prefix cache vLLM already
has.

\section{2. Related work}\label{related-work}

\textbf{Prefix caching.} vLLM's PagedAttention {[}1{]} shares fixed-size
KV blocks across requests with identical prefixes, and SGLang's
RadixAttention {[}8{]} keeps cached prefixes in a radix tree. Both are
exact, and both miss when reused text follows different preceding text,
which is the normal case for retrieved memories.

\textbf{Position-independent reuse.} Prompt Cache {[}9{]} precomputes
attention states for prompt modules at positions declared in a schema,
without correcting for the cross-module attention it skips. CacheGen
{[}10{]} compresses and streams KV so it can be loaded instead of
recomputed, which is complementary to our work.

\textbf{Correcting reused KV.} EPIC {[}3{]} calls the setting
position-independent caching. It observes that a chunk cached alone
treats its first tokens as an attention sink {[}11{]}, which is wrong
once the chunk sits mid-prompt, so it recomputes the first few tokens
(at most 32) of every chunk after the first. CacheBlend {[}4{]} finds
that about 10-15\% of tokens deviate strongly from their fresh KV, picks
them at an early layer and recomputes only those. On RAG benchmarks it
reports 2.2-3.3× lower TTFT than full prefill with quality within 0.02
F1. AgentKVShift {[}5{]} targets agent memory. Its spectral analysis
finds that most of a chunk's reuse error is one shared offset, so it
recomputes a few probe tokens, estimates the per-layer mean shift and
adds it to every other token. On LoCoMo with Mistral-7B-Instruct-v0.3
and the LiCoMemory memory system it reports 0.491 F1 against 0.509 for
full recompute at a 10\% budget (a 3.5\% relative gap), and it reports
2-3.5× prefill speedups on a single A100. It is our main comparison
point.

\textbf{Serving and agents.} LMCache {[}6{]}, from the same group as
CacheBlend, is the reference implementation of vLLM's KV-connector
interface and the open-source vLLM integration of non-prefix reuse we
build against. LoCoMo {[}7{]} provides very long multi-session
conversations with QA annotations. MemGPT {[}12{]} and A-Mem {[}13{]}
are memory systems that produce the retrieve-then-prompt pattern we
target. KVCOMM {[}14{]} and KVCMAS {[}15{]} share KV across agents and
matter only for our stretch goal.

\section{3. Proposed work}\label{proposed-work}

\subsection{3.1 Approach}\label{approach}

Each memory entry is a \emph{chunk}. The first time a chunk is retrieved
we prefill it alone and store its KV under a hash of its tokens. Later
prompts containing it load that KV, re-rotate the keys to the new
position and apply one of four corrections, which differ only in which
tokens they recompute:

\begin{longtable}[]{@{}
  >{\raggedright\arraybackslash}p{(\linewidth - 6\tabcolsep) * \real{0.2500}}
  >{\raggedright\arraybackslash}p{(\linewidth - 6\tabcolsep) * \real{0.2500}}
  >{\raggedright\arraybackslash}p{(\linewidth - 6\tabcolsep) * \real{0.2500}}
  >{\raggedright\arraybackslash}p{(\linewidth - 6\tabcolsep) * \real{0.2500}}@{}}
\toprule\noalign{}
\begin{minipage}[b]{\linewidth}\raggedright
\#
\end{minipage} & \begin{minipage}[b]{\linewidth}\raggedright
Method
\end{minipage} & \begin{minipage}[b]{\linewidth}\raggedright
Recomputes
\end{minipage} & \begin{minipage}[b]{\linewidth}\raggedright
Assumes the error is
\end{minipage} \\
\midrule\noalign{}
\endhead
\bottomrule\noalign{}
\endlastfoot
1 & Re-positioning only (baseline) & nothing & negligible \\
2 & Boundary (EPIC) & first 16-32 tokens of each chunk & at chunk
starts \\
3 & Dynamic selection (CacheBlend) & 10-15\% of tokens deviating most at
an early layer & sparse, anywhere \\
4 & Offset correction (AgentKVShift) & a few probes; their mean shift is
added to the rest & one shared shift per chunk \\
\end{longtable}

Every prompt is the system prompt, then the top-\emph{k} retrieved
entries, then the question, which is always computed fresh.

\subsection{3.2 Architecture}\label{architecture}

\begin{verbatim}
 LoCoMo conversations                         Synthetic Zipf pool
         |                                            |
 +-------v-----------------+                +---------v----------+
 | C1 Memory pipeline      |                | C4 Request gen     |
 | chunk, BM25/embed top-k |                | (tunable repeats)  |
 | -> retrieval.json       |                +---------+----------+
 +-------+-----------------+                          |
         | same saved retrievals for every config     |
    +----+----------------------+                     |
    v                           v                     v
 +----------------------+   +------------------------------------------+
 | C2 Transformers      |   | vLLM (unmodified, pinned)                |
 | chunk KV store       |   |   C3 KV connector: look up chunk KV,     |
 | RoPE re-rotation     |   |   re-rotate, concatenate as prefix;      |
 | per-layer recompute  |   |   vLLM prefills only the question        |
 | methods 1-4          |   +--------------------+---------------------+
 +----------+-----------+                        |
            v                                    v
   E0-E2: F1, prefill cost        C5 benchmark harness, E3-E5: TTFT, throughput
\end{verbatim}

\textbf{C1, memory pipeline.} Splits each LoCoMo conversation into
entries of a few hundred tokens, retrieves the top-\emph{k} per question
with BM25 or embeddings, and saves the results to disk. Every
configuration reads the same file, so no difference between methods can
come from retrieval.

\textbf{C2, Transformers implementation (the core).} A chunk KV store
keyed by token hash, RoPE re-rotation, and a custom forward pass that
steps through the decoder layers and, at each layer, recomputes a token
set \emph{S} against the assembled cache and writes the fresh KV back.
Each method is a function that returns \emph{S} (plus, for method 4, an
offset for the other tokens).

\textbf{C3, vLLM connector plugin.} vLLM can load an external KV
connector that supplies cached KV for the first \emph{N} tokens of a
prompt. Ours builds that prefix from the re-rotated KV of the retrieved
chunks, so vLLM computes only the question. This serves method 1 without
touching vLLM's source. Methods 2 to 4 recompute tokens inside the
prefix at every layer, which would need scheduler changes, so they stay
in Transformers.

\textbf{C4, synthetic workload} issues requests over a pool of entries
with Zipf-distributed popularity; the exponent sets how often entries
repeat. \textbf{C5, benchmark harness} drives vLLM's benchmark tool with
a custom dataset and records TTFT, latency and throughput.

\subsection{3.3 Novelty and scope}\label{novelty-and-scope}

The correction algorithms come from prior work, and we do not propose a
new one. What we add: a like-for-like comparison (published results use
different models, datasets and memory formats, while we run all four
methods with one model, one retrieval output and one budget accounting);
a new setting for EPIC and CacheBlend, which were evaluated on RAG
passages rather than conversational memory; non-prefix reuse served
through vLLM's public plugin interface instead of engine changes; and a
break-even analysis of how often memories must repeat before reuse beats
prefix caching, which the papers above do not report.

The core is C1, C2 and experiments E0-E2, which answer the research
question (quality against prefill cost) on their own. They measure
prefill time in Transformers, not serving latency or throughput. Those
belong to the secondary question, which C3-C5 and E3-E5 answer, and
that work starts only after the E0 correctness gate passes. If it does
not finish, the report makes no serving claims. A multi-agent workload, where several agents
read overlapping memories through one connector, is a stretch goal.
Changing the source of vLLM, Transformers or the model is out of scope.

\section{4. Plan of action}\label{plan-of-action}

\subsection{4.1 Resources}\label{resources}

\begin{itemize}
\tightlist
\item
  \textbf{Hardware:} one A100-80GB or H100 per run on PACE ICE, through
  the course allocation.
\item
  \textbf{Model and data:} Mistral-7B-Instruct-v0.3 {[}16{]} from
  HuggingFace, the model AgentKVShift reports on LoCoMo; LoCoMo from its
  public release (10 conversations, 1,986 questions).
\item
  \textbf{Software:} PyTorch, Transformers, \texttt{rank\_bm25},
  sentence-transformers, and vLLM pinned to 0.27.1 in a lock file. Code, configs and result logs live in a team GitHub repository.
\end{itemize}

The model's bf16 weights need about 15 GB of GPU memory. Its KV takes
128 KB per token (32 layers × 8 KV heads × 128 dims × K and V × 2
bytes), and LoCoMo conversations average about 16K tokens {[}7{]}, so
the KV for all ten conversations is about 20 GB and the chunk store fits
in GPU or host memory without distributed storage.

\subsection{4.2 Weekly schedule}\label{weekly-schedule}

Each component has an owner, and everyone reviews the E0 gate and the
final numbers.

\begin{longtable}[]{@{}
  >{\raggedright\arraybackslash}p{(\linewidth - 4\tabcolsep) * \real{0.3333}}
  >{\raggedright\arraybackslash}p{(\linewidth - 4\tabcolsep) * \real{0.3333}}
  >{\raggedright\arraybackslash}p{(\linewidth - 4\tabcolsep) * \real{0.3333}}@{}}
\toprule\noalign{}
\begin{minipage}[b]{\linewidth}\raggedright
Week
\end{minipage} & \begin{minipage}[b]{\linewidth}\raggedright
Work
\end{minipage} & \begin{minipage}[b]{\linewidth}\raggedright
Milestone
\end{minipage} \\
\midrule\noalign{}
\endhead
\bottomrule\noalign{}
\endlastfoot
6 (Sep 28-Oct 2) & Finalize scope, submit proposal & Proposal in \\
7 (Oct 5-9) & PACE setup, model and data download, full-prefill F1
harness (all) & Baseline F1 \\
8 (Oct 12-16) & C1 and saved retrievals (Prakhar); KV store, re-rotation
(Atharva); recompute loop (Teppei, Vishal) & Retrievals frozen \\
9 (Oct 19-23) & E0 tests; methods 1 and 2 (Atharva, Teppei) & E0 passes
(gate) \\
10 (Oct 26-30) & Methods 3 and 4 (Teppei, Vishal); first E1 runs
(Prakhar); study connector API (Atharva) & All methods run \\
11 (Nov 2-6) & Full E1 sweep and E2 (Prakhar, Teppei, Vishal); connector
prototype (Atharva) & Core E0-E2 done \\
12 (Nov 9-13) & Finish C3 (Atharva); C4 and C5 (Prakhar); E5 (Teppei);
report outline (Vishal) & E5 passes \\
13 (Nov 16-20) & E3 and E4 (Atharva, Prakhar); ablations (Teppei,
Vishal); draft report sections & All data in \\
14 (Nov 23-27) & Thanksgiving week: full report draft and figures;
buffer for slipped runs & Report drafted \\
15 (Nov 30-Dec 1) & Final edits, code cleanup, reproducibility README &
Submit Dec 1 \\
\end{longtable}

Experiments end on Nov 20 so the last week and a half go to writing. If
the core slips, E4 is cut first. The stretch goal starts only if E3-E5
finish early. E0-E2 are never cut.

\subsection{4.3 Risks}\label{risks}

vLLM's connector interface changes between releases and may make it
harder than expected to inject assembled KV; we pin one version and
follow LMCache's connector, and the core does not depend on vLLM. Reuse
may lose more quality than the papers report: AgentKVShift's own LoCoMo
table puts CacheBlend at 0.290 F1 against 0.509 for full recompute
(Mistral-7B, LiCoMemory, 10\% budget), and if we see the same we report
it. Finally, AgentKVShift runs on memory systems (A-Mem, LiCoMemory)
with generated metadata while our memories are plain conversation
chunks, so we report both absolute and relative F1 gaps.

\section{5. Evaluation and testing
method}\label{evaluation-and-testing-method}

\textbf{Setup.} Mistral-7B-Instruct-v0.3 on one PACE GPU, greedy
decoding, the same saved retrievals for every configuration, one warm-up
run and the median of three repetitions. Quality is LoCoMo F1, overall
and per question category, with full prefill of the same prompt as the
reference. Systems runs use the Zipf workload. Every result is stored
with its configuration (vLLM version, GPU, \emph{k}, budget, seed).

\begin{longtable}[]{@{}
  >{\raggedright\arraybackslash}p{(\linewidth - 6\tabcolsep) * \real{0.2500}}
  >{\raggedright\arraybackslash}p{(\linewidth - 6\tabcolsep) * \real{0.2500}}
  >{\raggedright\arraybackslash}p{(\linewidth - 6\tabcolsep) * \real{0.2500}}
  >{\raggedright\arraybackslash}p{(\linewidth - 6\tabcolsep) * \real{0.2500}}@{}}
\toprule\noalign{}
\begin{minipage}[b]{\linewidth}\raggedright
ID
\end{minipage} & \begin{minipage}[b]{\linewidth}\raggedright
Question
\end{minipage} & \begin{minipage}[b]{\linewidth}\raggedright
Stack
\end{minipage} & \begin{minipage}[b]{\linewidth}\raggedright
Measurement
\end{minipage} \\
\midrule\noalign{}
\endhead
\bottomrule\noalign{}
\endlastfoot
E0 & Is the implementation correct? & Transformers & (a) RoPE round
trip: keys shifted away and back equal the originals; (b) first layer:
shifted keys equal keys computed directly at the target position; (c)
full path: 100\% recompute equals full-prefill logits \\
E1 & How much quality does each method keep? & Transformers & LoCoMo F1
vs.~recompute budget (0-60\%), one line per method \\
E2 & How much prefill is removed? & Transformers & Prefill time and
tokens computed vs.~\emph{k} \\
E3 & Does the connector cut latency? & vLLM & TTFT vs.~\emph{k}: full
prefill, prefix caching, connector \\
E4 & When does reuse beat prefix caching? & vLLM, Zipf & Throughput and
TTFT as reuse rate goes from \textasciitilde0\% to
\textasciitilde100\% \\
E5 & Does the connector match the reference? & vLLM vs.~HF & Method 1:
identical tokenization and matching sampled KV tensors; then
near-identical greedy outputs and F1 as end-to-end checks \\
\end{longtable}

\textbf{Validation.} E0 is a hard gate: if any of its three checks
fails, every later number is meaningless, so nothing else is reported
until all three pass. E0 tests only cases where an exact match is
expected. Once a chunk's preceding context changes, re-rotation alone
cannot match full prefill, and that gap is what E1 measures. Unit tests cover RoPE round
trips, hash collisions in the chunk store and cache assembly for
out-of-order chunks. E5 ties the vLLM path to the Transformers path, and
we compare our full-prefill and method 4 F1 with the numbers in {[}5{]}
and explain any gap. Recomputed tokens are counted the same way for
every method, including tokens touched at a selection layer. If time
permits we ablate K-only versus K-and-V offsets in method 4, random
versus top-divergence probes, and per-category F1, since both papers
suggest multi-hop questions suffer most.

\textbf{Success criteria.} (1) AgentKVShift stays within 0.02 absolute
F1 of full prefill at a 10-15\% recompute budget, while re-positioning
only scores clearly lower. (2) For CacheBlend and EPIC we report the
full quality-cost curve up to 60\% and the smallest budget, if any, that
comes within 0.02 F1. We set no fixed threshold for them, because {[}5{]}
reports CacheBlend at 0.290 F1 against 0.509 at a 10\% budget and notes
that earlier methods may need 45-55\% recompute to approach full
quality. (3) AgentKVShift cuts prefill time by at least 2× when 6 or
more entries are retrieved. (4) For the secondary question, the
connector beats both full prefill and prefix caching on TTFT when
entries repeat in different orders. Missing a criterion is still a
result, and the report will state what we measured.

\section{6. Bibliography}\label{bibliography}

\small

{[}1{]} W. Kwon, Z. Li, S. Zhuang, Y. Sheng, L. Zheng, C. H. Yu, J. E.
Gonzalez, H. Zhang, I. Stoica. Efficient Memory Management for Large
Language Model Serving with PagedAttention. \emph{SOSP}, 2023.
arXiv:2309.06180. https://github.com/vllm-project/vllm

{[}2{]} J. Su, Y. Lu, S. Pan, A. Murtadha, B. Wen, Y. Liu. RoFormer:
Enhanced Transformer with Rotary Position Embedding. arXiv:2104.09864,
2021.

{[}3{]} J. Hu, W. Huang, W. Wang, H. Wang, T. Hu, Q. Zhang, H. Feng, X.
Chen, Y. Shan, T. Xie. EPIC: Efficient Position-Independent Caching for
Serving Large Language Models. \emph{ICML}, 2025. arXiv:2410.15332.

{[}4{]} J. Yao, H. Li, Y. Liu, S. Ray, Y. Cheng, Q. Zhang, K. Du, S. Lu,
J. Jiang. CacheBlend: Fast Large Language Model Serving for RAG with
Cached Knowledge Fusion. \emph{EuroSys}, 2025. arXiv:2405.16444.

{[}5{]} N. P. Pandey, J. Kong, L. Hu, Q. Zhao, Y. Zhao, O. Gungor, H.
Zhang, T. Rosing. AgentKVShift: Efficient KV Cache Reuse for Agentic
Memory Systems. arXiv:2607.21604, 2026.

{[}6{]} LMCache. https://github.com/LMCache/LMCache (issues \#4476 and
\#5101 on the blend path).

{[}7{]} A. Maharana, D.-H. Lee, S. Tulyakov, M. Bansal, F. Barbieri, Y.
Fang. Evaluating Very Long-Term Conversational Memory of LLM Agents.
\emph{ACL}, 2024. arXiv:2402.17753.
https://github.com/snap-research/locomo

{[}8{]} L. Zheng et al.~SGLang: Efficient Execution of Structured
Language Model Programs. \emph{NeurIPS}, 2024. arXiv:2312.07104.

{[}9{]} I. Gim, G. Chen, S. Lee, N. Sarda, A. Khandelwal, L. Zhong.
Prompt Cache: Modular Attention Reuse for Low-Latency Inference.
\emph{MLSys}, 2024. arXiv:2311.04934.

{[}10{]} Y. Liu et al.~CacheGen: KV Cache Compression and Streaming for
Fast Large Language Model Serving. \emph{SIGCOMM}, 2024.
arXiv:2310.07240.

{[}11{]} G. Xiao, Y. Tian, B. Chen, S. Han, M. Lewis. Efficient
Streaming Language Models with Attention Sinks. \emph{ICLR}, 2024.
arXiv:2309.17453.

{[}12{]} C. Packer et al.~MemGPT: Towards LLMs as Operating Systems.
arXiv:2310.08560, 2023.

{[}13{]} W. Xu, Z. Liang, K. Mei, H. Gao, J. Tan, Y. Zhang. A-Mem:
Agentic Memory for LLM Agents. \emph{NeurIPS}, 2025. arXiv:2502.12110.

{[}14{]} H. Ye et al.~KVCOMM: Online Cross-context KV-cache
Communication for Efficient LLM-based Multi-agent Systems.
\emph{NeurIPS}, 2025. arXiv:2510.12872.
https://github.com/FastMAS/KVCOMM

{[}15{]} H. Jeon, H. Ha, S. Lee, B. Kang, J.-J. Kim. KVCMAS: Efficient
KV Cache Correction for Shared Context in Multi-Agent Systems.
arXiv:2609.34060, 2026.

{[}16{]} A. Q. Jiang et al.~Mistral 7B. arXiv:2310.06825, 2023.

\end{document}
